\documentclass[12pt]{article}
\usepackage[utf8]{inputenc}
\usepackage[T1]{fontenc}
\usepackage{hyperref}
\title{Wyandot Language Resources for K-12 Curriculum Development}
\author{Compiled by Claude Sonnet 5.5 (claude-sonnet-5-5) for Andr\'as Kornai}
\date{}
\begin{document}
\maketitle

\section{Introduction}
Wyandot (ISO 639-3: wyn) is an Iroquoian language, dormant as a first language and used by some younger people as a second language. The Wyandotte Nation (Oklahoma) runs language work, and Wendat communities in Quebec also offer classes. This survey is short and based on a few web searches. Little material was found that is ready for classroom use, and several leads below could not be confirmed on primary sites.

\section{Primary Educational Resources}
\begin{itemize}
\item Wikipedia reports that the Wyandotte Nation teaches Wyandot to K--4 students in Wyandotte Public Schools and at its preschool, and has made online self-study lessons: \url{https://en.wikipedia.org/wiki/Wyandot_language}. This is a secondary source, and the lessons themselves were not located.
\item \emph{DSMW -- A Wandat (Wyandot) Language Workbook}, listed in the ONLC bookstore: \url{https://www.onlc.ca/bookstore/p/dsmw-a-wandat-wyandot-language-workbook}. I saw only the listing, not the contents.
\end{itemize}

\section{Linguistic and Reference Materials}
\begin{itemize}
\item John Steckley, tribal linguist for the Wyandotte Nation, published the first Huron-English dictionary in 2007: \url{https://www.nccie.ca/story/john-steckley-wyandot-tribal-linguist/}. The main source texts are about forty narratives that Marius Barbeau recorded in 1911--1912.
\item Related Wendat materials from 17th and 18th century Jesuit missionaries, compiled by Craig Kopris, are mentioned in the same source. No separate URL was seen.
\item An Illinois State University dissertation on reclaiming Wyandot, done with the Wyandot(te) community: \url{https://ir.library.illinoisstate.edu/etd/2297}.
\item Omniglot overview of Wyandot: \url{https://mail.omniglot.com/writing/wyandot.htm}.
\end{itemize}

\section{Religious and Community Resources}
\begin{itemize}
\item The Wyandotte Nation culture page mentions classes at the Cultural Center: \url{https://wyandotte-nation.org/culture/}. I could not confirm a dedicated language department page.
\end{itemize}

\section{Digital and Audio Resources}
No audio or video resources for Wyandot were found in these searches.

\section{Curriculum Development Context}
Wyandot has a small number of second-language learners and a few existing school and preschool classes. Curriculum building would rely on the Steckley dictionary and the Barbeau-derived texts. Anyone planning this should contact the Wyandotte Nation language program directly, since I could not verify what lessons or materials are currently available.

\end{document}